% font: footnotesize
% TAB-14: generated by harness/analysis/paper_outputs.py; do not edit (every cell's source is in TAB-14.cells.json)
\begin{tabularx}{\linewidth}{>{\hsize=2.048\hsize\raggedright\hyphenpenalty=10000\exhyphenpenalty=10000\arraybackslash}X>{\hsize=0.517\hsize\raggedright\hyphenpenalty=10000\exhyphenpenalty=10000\arraybackslash}X>{\hsize=0.408\hsize\raggedleft\hyphenpenalty=10000\exhyphenpenalty=10000\arraybackslash}X>{\hsize=0.576\hsize\raggedleft\hyphenpenalty=10000\exhyphenpenalty=10000\arraybackslash}X>{\hsize=0.439\hsize\raggedleft\hyphenpenalty=10000\exhyphenpenalty=10000\arraybackslash}X>{\hsize=2.013\hsize\raggedright\hyphenpenalty=10000\exhyphenpenalty=10000\arraybackslash}X}
\toprule
Setting & Shape & List cost per puzzle (\$) & Extra list cost of the picker call alone (\$) & List cost per correct answer (\$) & Label \\
\midrule
Haiku team & chains (1\ensuremath{\times}6) & 0.00623 &  & 0.0119\phantom{0} & exploratory (no Holm, no verdict) \\
Haiku team & pieces (6\ensuremath{\times}1) & 0.00649 &  & 0.00895 & exploratory (no Holm, no verdict) \\
Haiku team & pooled & 0.00636 &  &  & exploratory (no Holm, no verdict) \\
\addlinespace
Haiku picker over 8 samples & chains (1\ensuremath{\times}6) & 0.0103\phantom{0} & 0.00154 & 0.0165\phantom{0} & exploratory (no Holm, no verdict) \\
Haiku picker over 8 samples & pieces (6\ensuremath{\times}1) & 0.0101\phantom{0} & 0.00142 & 0.0150\phantom{0} & exploratory (no Holm, no verdict) \\
Haiku picker over 8 samples & pooled & 0.0102\phantom{0} & 0.00148 &  & exploratory (no Holm, no verdict) \\
\addlinespace
Step voting on pieces & pieces (6\ensuremath{\times}1) & 0.00351 &  & 0.00684 & unadjusted, exploratory, outside the pre-registered family \\
\addlinespace
One Haiku-medium sample (the first 8 samples) & chains (1\ensuremath{\times}6) &  &  &  & comparator; exploratory (no Holm, no verdict) \\
One Haiku-medium sample (the first 8 samples) & pieces (6\ensuremath{\times}1) &  &  &  & comparator; exploratory (no Holm, no verdict) \\
One Haiku-medium sample (the first 8 samples) & pooled &  &  &  & comparator; exploratory (no Holm, no verdict) \\
\addlinespace
Haiku-medium vote of the first 8 samples & chains (1\ensuremath{\times}6) & 0.00879 &  &  & comparator; exploratory (no Holm, no verdict) \\
Haiku-medium vote of the first 8 samples & pieces (6\ensuremath{\times}1) & 0.00873 &  &  & comparator; exploratory (no Holm, no verdict) \\
Haiku-medium vote of the first 8 samples & pooled & 0.00876 &  &  & comparator; exploratory (no Holm, no verdict) \\
\addlinespace
Any of the first 8 Haiku-medium samples right (the ceiling for a picker) & chains (1\ensuremath{\times}6) &  &  &  & comparator; exploratory (no Holm, no verdict) \\
Any of the first 8 Haiku-medium samples right (the ceiling for a picker) & pieces (6\ensuremath{\times}1) &  &  &  & comparator; exploratory (no Holm, no verdict) \\
Any of the first 8 Haiku-medium samples right (the ceiling for a picker) & pooled &  &  &  & comparator; exploratory (no Holm, no verdict) \\
\addlinespace
Opus-medium, main run & chains (1\ensuremath{\times}6) & 0.0402\phantom{0} &  & 0.0418\phantom{0} & comparator; pre-specified descriptive (Study~1, as ⟦TAB-04⟧ prints it) \\
Opus-medium, main run & pieces (6\ensuremath{\times}1) & 0.0481\phantom{0} &  & 0.0481\phantom{0} & comparator; pre-specified descriptive (Study~1, as ⟦TAB-04⟧ prints it) \\
Opus-medium, main run & pooled &  &  &  & comparator; exploratory (no Holm, no verdict) \\
\addlinespace
Step voting, main run & chains (1\ensuremath{\times}6) & 0.00356 &  & 0.00662 & comparator; pre-specified descriptive (Study~1, as ⟦TAB-04⟧ prints it) \\
\bottomrule
\end{tabularx}
% note: On pieces the Haiku team is compared with step voting on pieces, both in the top-up run (across runs: no); in the main run, step voting ran on chains only.
% note: The top-up is exploratory (no Holm, no verdict). Step voting on pieces against one Opus-medium answer from the main run is the pieces counterpart of Test~4, not one of Tests~1--5: it is labelled unadjusted (exploratory, outside Holm) and is outside the pre-registered family; the drift check between runs: not fired.
% note: A comparison across runs sets a top-up setting against a main-run setting run at another time, and carries the status of the drift check between runs. The drift check, set by the top-up run's specification (frozen with Lock~4), re-ran Haiku-medium at the main run's Haiku-medium settings on 20 puzzles in both shapes (40 puzzle-shape pairs, 5 samples each, 200 calls) and set it against the main run's first 5 Haiku-medium samples on each of the same puzzles: pooled $d$ (top-up minus main run)~=~+1.50 points (95\% interval \ensuremath{-}8.00 to +10.00); mean output tokens per call changed by +0.29\%; status: not fired (a flag, not a gate; these answers are never mixed into Study~1).
% note: Rule for the drift check between runs: fires if the pooled 95\% interval excludes 0 and \ensuremath{|}$d$\ensuremath{|} \ensuremath{>}~10 points, or mean output tokens move by more than 25\%; a flag, not a gate.
% note: Picker cost: the list cost per puzzle is the total, the 8 main-run samples plus the picker call; the extra list cost is the picker call alone, on top of the 8 samples that the vote of 8 also uses. Cost per correct answer is given for the total only.
% note: A Label cell that starts with ``comparator'' marks a main-run comparator. The rows of the first 8 Haiku-medium samples and the pooled Opus-medium row are rebuilt by the top-up analysis from the main run's log; the Opus-medium rows per shape and step voting on chains print Study~1's own stamped values (the rebuild gives the same accuracies). The vote of 8 is priced at the list cost of its 8 samples, the picker's stamped total minus its stamped extra. The one-sample row has no stamped cost. The any-of row is a ceiling that needs the answer key, not a setup one could run, so it is given no cost.
% note: Format compliance is the share of answered replies whose format was ok, where a file stamps it: the Haiku team's (top-up file) and Opus-medium's (Study~1). Readable picks is the share of picker replies that named one of the samples. Other empty cells are not stamped.
% note: ``(exact)'' marks an exact (Clopper--Pearson) 95\% interval. The main-run rows that Study~1 also reports print the interval Study~1 stamps for the same result (as ⟦TAB-04⟧ prints it), so the paper gives one interval per result: Opus-medium's one sample on chains and on pieces (exact) and step voting on chains (Study~1's cluster bootstrap). The top-up file's own bootstraps for these rows are not printed; for Opus-medium on pieces it collapsed to one value. The same exact interval for the Opus-medium pieces cell is stamped by free analysis 4 (post hoc): \mbox{[95.49, 100.00]}. Every other row keeps the top-up file's cluster bootstrap; Study~1's pre-registered settings table (⟦TAB-04⟧) has no row for the top-up designs, the first 8 Haiku-medium samples or pooled shapes.
